% Research candidate - do NOT import into main.tex until proof audit / Lean review.
% The finite identities are elementary; the uniform dynamical bounds remain open.
\subsection*{Return-occupation potential and an exact second-moment identity}
Let $T:\mathbb N\to\mathbb N$ be a deterministic map, $a\ge0$ a
first-passage threshold, and $I\subset(a,\infty)\cap\mathbb N$ finite.
Fix a finitely supported probability law $\mu$ on starting points for which
$\tau_a(n)=\inf\{t\ge0:T^t(n)\le a\}<\infty$ holds $\mu$-almost surely.
Set
\[
  N_I(n)=\#\{0\le t<\tau_a(n):T^t(n)\in I\},\qquad
  h_{a,\mu}(y)=\mu\{n:\exists\,t<\tau_a(n),\ T^t(n)=y\}.
\]
For each $y$ with $h_{a,\mu}(y)>0$, define $r_{I,a}(y)$ as the number of
\emph{later} visits to $I$ along the orbit starting at $y$ before its
first passage through $a$.

\paragraph{Proposition (Exact killed-potential moment formula).}
Under the preceding assumptions,
\begin{equation}\label{eq:killed-return-moment}
  \boxed{\displaystyle
  \mathbb E_\mu N_I^2
     =\sum_{\substack{y\in I\\ h_{a,\mu}(y)>0}}
       h_{a,\mu}(y)\bigl(1+2r_{I,a}(y)\bigr).}
\end{equation}
\emph{Proof.} A deterministic orbit which reaches $\le a$ cannot repeat
any state before the passage, as a repetition would trap it in a periodic
orbit above $a$. If its visits to $I$ are $y_1,\ldots,y_N$ in temporal
order, then $r_{I,a}(y_i)=N-i$. Consequently
\[
  N^2=\sum_{i=1}^{N}\bigl(1+2(N-i)\bigr).
\]
Multiply by the starting-point weight and interchange the two finite
sums. Since each $y\in I$ is visited at most once, the coefficient of
$1+2r_{I,a}(y)$ is precisely $h_{a,\mu}(y)$.\hfill$\square$

\paragraph{Equivalent Green-operator formulation.}
Let $q_1$ denote the first-entrance distribution on $I$, and let $Q_I$
be the deterministic partial first-return push-forward before passage
through $a$. On the reachable return graph, which is acyclic under the
stated finite-passage assumption, put
\[
  \nu_I=\sum_{j\ge0}Q_I^j q_1,\qquad
  \sigma_I=\sum_{j\ge0}(j+1)Q_I^j q_1.
\]
These sums are finite on the support reachable from $q_1$. Exactly,
\begin{equation}\label{eq:double-green}
 \boxed{\displaystyle
 \nu_I-Q_I\nu_I=q_1,\quad
 \sigma_I-Q_I\sigma_I=\nu_I,\quad
 \mathbb E_\mu N_I^2=2\sigma_I(I)-\nu_I(I).}
\end{equation}
In particular, $\nu_I(y)=h_{a,\mu}(y)$: for finite-passage trajectories,
the occupation measure is a \emph{killed predecessor potential} and not
an unconditioned stationary or iid entrance law.

\paragraph{A precise limitation of naive potential weights.}
The first Green identity implies, at every shell state with
$q_1(y)=0$ and $\nu_I(y)>0$, that $(Q_I\nu_I)(y)=\nu_I(y)$.
Consequently, $\nu_I$ \emph{cannot} satisfy a strict pointwise
sub-eigenmeasure inequality $Q_I\nu_I\le\rho\nu_I$ for any $\rho<1$.
The occupation potential alone is therefore not the hoped-for strict
contraction weight. There is an exact, but quantitatively nontrivial,
remedy: for any $0<\rho<1$, define
\[
 \nu_{I,\rho}=\sum_{j\ge0}\rho^{-j}Q_I^j q_1.
 \quad\text{Then}\quad
 Q_I\nu_{I,\rho}=\rho(\nu_{I,\rho}-q_1)\le\rho\nu_{I,\rho}.
\]
However, $\nu_{I,\rho}(I)=\mathbb E_\mu\sum_{j=0}^{N_I-1}\rho^{-j}$;
bounding this uniformly in the shell scale is an \emph{exponential-moment}
problem, strictly stronger than the second-moment bound we seek.

\paragraph{Relation to the manuscript's predecessor potential.}
If $I$ contains \emph{no periodic point}, define $r_I(y)$ by counting
all later visits to $I$ in the forward orbit of $y$, without a passage
cutoff. Every such orbit can visit any $y\in I$ at most once. Suppose the
logarithmic densities $h(y)=\delta_{\log}\{n:\exists t,\ T^t(n)=y\}$
exist for $y\in I$. Then, for the distinct-state full-orbit visit count
$N_I^{\mathrm{full}}(n)$,
\begin{equation}\label{eq:full-log-second-moment}
 \boxed{\displaystyle
 \lim_{M\to\infty}\frac1{H_M}
   \sum_{n=1}^{M}\frac{N_I^{\mathrm{full}}(n)^2}{n}
 =\sum_{y\in I}h(y)\bigl(1+2r_I(y)\bigr),\quad
 H_M=\sum_{n\le M}\frac1n.}
\end{equation}
Indeed, apply the same pointwise identity and pass through the \emph{finite}
outer sum. No separate assumption on joint predecessor-event densities
is needed. This does \emph{not} imply a bound uniform in the moving
shells $I=I_{X,j}$. Nor can one replace $h_{a,\mu}$ by $h$ in
\eqref{eq:killed-return-moment}: first-passage killing matters.

\paragraph{Open quantitative target.}
For logarithmic entrance laws, \emph{if} the killed hitting-event densities
$h_{a_X}(y)$ exist, the desired uniform second moments are equivalent to
\[
 \sup_{X,j}\sum_{\substack{y\in I_{X,j}\\h_{a_X}(y)>0}}
 h_{a_X}(y)\bigl(1+2r_{I_{X,j},a_X}(y)\bigr)<\infty.
\]
Neither existence of all these killed densities nor the uniform bound
is established here. The exact identities themselves require neither a
Markov approximation nor independence of Syracuse valuations.
